\section{方法}
\label{sec:method-zh}

\subsection{任务定义}

系统根据条件组合 $c$ 生成量化的符号乐谱片段 $x$。条件可包括音级集合 $S$、六维音程向量 $v(S)$、十二音序列 $r$、序列形式 $f\in\{P_n,R_n,I_n,RI_n\}$、节奏轮廓、姿态标签、乐器、声部数和小节数。请求输出包含 4--16 小节和 2--8 个声部，并以事件标记、MusicXML 和乐谱层面分析报告的形式返回。图~\ref{fig:method-pipeline-zh} 概括了该流程。系统不使用波形或其他音频表示。

条件语言模型分解为
\begin{equation}
  p_\theta(x\mid c)=\prod_{t=1}^{T}p_\theta(x_t\mid c,x_{<t}),
  \label{eq:conditional-lm-zh}
\end{equation}
其中 $x_t$ 为乐谱事件标记。研究目标是在下一标记建模的基础上加入显式且可检查的后调性约束。

\subsection{合法的程序化语料}

由于大量 1945 年后的艺术音乐作品仍受版权保护，项目内置语料由程序生成，而非抓取已出版乐谱。每个样本指定 4--16 小节、2--8 个声部、一个乐器标签、七种节奏轮廓之一和七种姿态标签之一。非序列样本带有音级集合与音程向量；序列样本则带有有效十二音序列以及 48 种移调 $P/R/I/RI$ 形式之一。这种划分避免要求同一片段既限制在小音级集合内，又完成全部十二个音级这一不相容目标。

规则生成器为每个声部建立节奏骨架，实施与姿态相关的起音、时值、休止和音区操作，并在请求乐器音域内选择音高。时间量化为 0.25 quarterLength。每个目标中实际生成的随机密度曲线仅保存用于隐藏评估，不进入模型可见的条件元数据，也不参与候选排序。

\subsection{条件与事件表示}

条件被序列化为以 \texttt{BOS} 开始、以 \texttt{SEP} 结束的前缀。前缀可包含音级、音程向量、序列音级、序列形式、节奏轮廓、姿态、声部数、小节数和乐器标记。显式的 \texttt{NO\_PCSET}、\texttt{NO\_ROW}、\texttt{NO\_PROFILE} 与 \texttt{NO\_GESTURE} 符号表示对应条件缺失。

乐谱主体由 \texttt{BAR}、\texttt{TIME\_SHIFT}、\texttt{VOICE}、\texttt{PITCH}、\texttt{OCTAVE}、\texttt{DUR} 和 \texttt{REST} 事件组成，并以 \texttt{EOS} 结束。有限状态语法约束解码顺序、可用声部索引、乐器音域和请求小节数，解析器把解码事件限制在请求跨度内。

许多完整片段超过模型 256 标记的上下文。训练采用保留条件的窗口，而不是静默截断序列。每个窗口重复完整条件前缀，并在容量允许时保留前序乐谱上下文；损失仅作用于当前乐谱主体目标片段。十轮覆盖周期对目标窗口分区，使每个保存的主体标记在一个周期中恰好参与一次训练。验证和教师强制评估枚举全部窗口。报告的标记准确率与交叉熵只是主体标记诊断，不代表作曲质量。

\subsection{条件 Transformer}

生成器采用六层因果 Transformer \citep{vaswani2017attention}，隐藏维度为 384，包含六个注意力头、宽度为 $4d$ 的 GELU 前馈模块、可学习位置嵌入、预归一化和 0.1 的 dropout。训练目标是在未遮蔽的乐谱主体标签上计算教师强制负对数似然：
\begin{equation}
  \mathcal{L}_{\mathrm{NLL}}(\theta)
  =-\sum_{t\in\mathcal{T}}\log p_\theta(z_t\mid z_{<t}),
  \label{eq:nll-zh}
\end{equation}
其中 $\mathcal{T}$ 表示当前目标窗口。系统不使用预训练语言模型或外部乐谱语料。

\subsection{约束引导的候选排序}

推断时，拟议解码器采样 $K$ 个受语法约束的候选，并在乐谱层面对每个候选进行分析。候选选择定义为
\begin{equation}
  x^*=\arg\min_{x\in\mathcal{C}_K}\sum_j\lambda_j m_j(x,c),
  \label{eq:reranking-zh}
\end{equation}
其中 $m_j$ 包括音级集合覆盖缺失、集合外音级、音程向量距离、循环序列顺序误差、不完整的严格序列形式、不完整的十二音聚合、节奏轮廓距离、姿态不一致、乐器音域违反、内容跨度不足和请求声部缺失。若无音级集合目标，则音级集合项不启用；若无序列目标，则严格形式和聚合项不启用。隐藏的随机密度曲线不参与该目标。这些不可微项仅用于候选选择，不被表述为训练损失。

主要受控比较使用同一个已训练生成器：$K=1$ 为普通解码器，$K=4$ 为引导式解码器。每个测试条件使用确定性的随机数流，使 $K=4$ 的首个候选与对应的 $K=1$ 输出相同。条件移除模型使用相同语料样本和目标事件，只修改模型可见的前缀字段；评估仍使用原始隐藏目标。

\subsection{分析与 MusicXML 导出}

每个选中片段均计算音级集合覆盖率、精确率和 Jaccard 相似度、生成音程向量、音程向量距离、循环序列顺序准确率、严格完整形式准确率、十二音聚合完成率、密度曲线、节奏轮廓距离、隐藏密度曲线误差、姿态特征、姿态一致性、音区跨度、音域违反率、内容跨度遵循率以及请求声部遵循率。严格序列形式准确率按不重叠的完整十二音块计算，并与循环序列顺序准确率相互独立。

MusicXML 导出器创建请求数量的乐谱声部和小节，补齐末尾空小节，并移除请求跨度之外的事件 \citep{musicxml2021spec}。系统分别记录结构解析成功率、请求小节数遵循率和请求声部数遵循率。
